% OVERLEAF WILL NOT BUILD THIS FILE. It \inputs docs/latex/coursework_preamble.tex by a
% relative path that climbs above the project root (so does reference_docs/cesc470_macros.tex),
% and an Overleaf project is self-contained: "File `../../../../docs/latex/coursework_preamble.tex' not found."
% Upload /home/devel/electrical_notes/content/cesc_470/hw/hw01/overleaf/p09_cpi_and_execution_time.tex instead -- docs/latex/flatten_tex.sh writes it.
% Edit THIS file, never that generated copy -- docs/latex/flatten_tex.sh overwrites it. Regenerate with:
%   docs/latex/flatten_tex.sh content/cesc_470/hw/hw01   (run from /home/devel/electrical_notes)
% ─────────────────────────────────────────────────────────────────────────────
% SHARED coursework preamble — every course inputs THIS.
%
% PREAMBLE ONLY. No \begin{document}, no course-specific notation.
%
% Course-specific macros live at COURSE level, in
%   content/<course>/reference_docs/<course>_macros.tex
% which is inputted AFTER this file. ONE per course, shared by hw/qz/exam --
% macros under a single kind are unreachable from a sibling kind.
%
% That split is the point: the page style, the answer box, and the problem
% header are identical everywhere, so they are defined once; \conv and \ustep
% mean something in DSP and nothing in a networks course, so they are not.
%
% A per-problem file therefore starts:
%   % ─────────────────────────────────────────────────────────────────────────────
% SHARED coursework preamble — every course inputs THIS.
%
% PREAMBLE ONLY. No \begin{document}, no course-specific notation.
%
% Course-specific macros live at COURSE level, in
%   content/<course>/reference_docs/<course>_macros.tex
% which is inputted AFTER this file. ONE per course, shared by hw/qz/exam --
% macros under a single kind are unreachable from a sibling kind.
%
% That split is the point: the page style, the answer box, and the problem
% header are identical everywhere, so they are defined once; \conv and \ustep
% mean something in DSP and nothing in a networks course, so they are not.
%
% A per-problem file therefore starts:
%   % ─────────────────────────────────────────────────────────────────────────────
% SHARED coursework preamble — every course inputs THIS.
%
% PREAMBLE ONLY. No \begin{document}, no course-specific notation.
%
% Course-specific macros live at COURSE level, in
%   content/<course>/reference_docs/<course>_macros.tex
% which is inputted AFTER this file. ONE per course, shared by hw/qz/exam --
% macros under a single kind are unreachable from a sibling kind.
%
% That split is the point: the page style, the answer box, and the problem
% header are identical everywhere, so they are defined once; \conv and \ustep
% mean something in DSP and nothing in a networks course, so they are not.
%
% A per-problem file therefore starts:
%   \input{../../../../docs/latex/coursework_preamble.tex}
%   \input{../../reference_docs/cesc470_macros.tex}
%   \begin{document}
%
% Build-check one file, or a whole assignment:
%   docs/latex/build_tex.sh content/cesc_470/hw/hw01/p01_five_components.tex
%   docs/latex/build_tex.sh content/cesc_470/hw/hw01
%
% Doctrine: docs/directives/coursework-solutions.md
% ─────────────────────────────────────────────────────────────────────────────

\documentclass[11pt]{article}

\usepackage[margin=1in]{geometry}
\usepackage{amsmath, amssymb, mathtools}
\usepackage{graphicx}
\usepackage{float}
\usepackage{booktabs}
\usepackage{longtable}
\usepackage{enumitem}
\usepackage{siunitx}
\usepackage[dvipsnames]{xcolor}
\usepackage{tcolorbox}
\usepackage{fancyhdr}
\usepackage{caption}
\usepackage{tikz}
\usepackage{pgfplots}
\pgfplotsset{compat=1.18}
\usepackage[colorlinks=true, allcolors=NavyBlue]{hyperref}

\captionsetup{font=small, labelfont=bf}
\setlist[enumerate]{itemsep=2pt, topsep=4pt}
\setlist[itemize]{itemsep=2pt, topsep=4pt}

% --- final-answer box -------------------------------------------------------
% Every problem file ends with its result in one of these. The condensed
% solutions document is assembled by lifting exactly these.
\newtcolorbox{answerbox}[1][]{
  colback=Green!6, colframe=Green!55!black,
  boxrule=0.6pt, arc=2pt, left=6pt, right=6pt, top=4pt, bottom=4pt,
  #1
}

% --- citation to course material --------------------------------------------
% \cite{Module 01, PDF p55}  ->  a small italic parenthetical.
% Solutions cite the SLIDE/PAGE a formula came from so a grader (or a future
% reader) can check the claim against what was actually taught. See the
% directive: every non-obvious step carries one.
\newcommand{\srcref}[1]{{\small\itshape(#1)}}

% --- a step that came from OUTSIDE the course materials ---------------------
% Used only where the course is genuinely silent and the problem asks for
% outside research. Marked visually so it is never mistaken for lecture content.
\newcommand{\extref}[1]{{\small\itshape\color{Sepia}[external: #1]}}

% --- callout for the trap in a problem --------------------------------------
\newtcolorbox{trap}[1][]{
  colback=BurntOrange!7, colframe=BurntOrange!70!black,
  boxrule=0.6pt, arc=2pt, left=6pt, right=6pt, top=4pt, bottom=4pt,
  fonttitle=\bfseries, title=Watch out, #1
}

% --- per-problem header -----------------------------------------------------
% \problemheader{8}{15 pts}{Target clock rate}
% The optional 4-argument form carries a learning outcome where the course
% states one (CESC 410 does; CESC 470's HW1 does not).
\newcommand{\problemheader}[3]{%
  \begin{center}\large\bfseries
    Problem #1 \quad\normalsize\mdseries(#2)\\[2pt]
    \large\bfseries #3
  \end{center}
  \vspace{2pt}\hrule\vspace{8pt}
}
\newcommand{\problemheaderlo}[4]{%
  \begin{center}\large\bfseries
    Problem #1 \quad\normalsize\mdseries(#2, #3)\\[2pt]
    \large\bfseries #4
  \end{center}
  \vspace{2pt}\hrule\vspace{8pt}
}

% Course name in the running head. Defined BEFORE \lhead references it, and
% \renewcommand'd by each course's macros file (inputted after this one).
\providecommand{\coursename}{}

\pagestyle{fancy}
\fancyhf{}
\lhead{\small\coursename}
\rhead{\small\thepage}
\renewcommand{\headrulewidth}{0.3pt}

%   % CESC 470 Computer Architecture — course-specific macros ONLY.
%
% Inputted AFTER docs/latex/coursework_preamble.tex, which carries everything
% shared (answerbox, problemheader, srcref, page style). Put nothing general
% here: if another course would want it, it belongs in the shared preamble.

\renewcommand{\coursename}{CESC 470}

% --- performance-equation notation ------------------------------------------
% The course writes these in words on the slides; these render them compactly
% and, more importantly, IDENTICALLY across every problem file.
\newcommand{\CPUtime}{T_{\mathrm{CPU}}}
\newcommand{\IC}{\mathrm{IC}}                    % instruction count
\newcommand{\CPI}{\mathrm{CPI}}
\newcommand{\cycletime}{t_{\mathrm{cycle}}}
\newcommand{\clockrate}{f_{\mathrm{clk}}}
\newcommand{\cycles}{N_{\mathrm{cycles}}}
\newcommand{\speedup}{S}

% Amdahl's law, in the form the slides state it (Module 01, PDF p86, slide 49).
\newcommand{\amdahl}[3]{#1 + \dfrac{#2}{#3}}

%   \begin{document}
%
% Build-check one file, or a whole assignment:
%   docs/latex/build_tex.sh content/cesc_470/hw/hw01/p01_five_components.tex
%   docs/latex/build_tex.sh content/cesc_470/hw/hw01
%
% Doctrine: docs/directives/coursework-solutions.md
% ─────────────────────────────────────────────────────────────────────────────

\documentclass[11pt]{article}

\usepackage[margin=1in]{geometry}
\usepackage{amsmath, amssymb, mathtools}
\usepackage{graphicx}
\usepackage{float}
\usepackage{booktabs}
\usepackage{longtable}
\usepackage{enumitem}
\usepackage{siunitx}
\usepackage[dvipsnames]{xcolor}
\usepackage{tcolorbox}
\usepackage{fancyhdr}
\usepackage{caption}
\usepackage{tikz}
\usepackage{pgfplots}
\pgfplotsset{compat=1.18}
\usepackage[colorlinks=true, allcolors=NavyBlue]{hyperref}

\captionsetup{font=small, labelfont=bf}
\setlist[enumerate]{itemsep=2pt, topsep=4pt}
\setlist[itemize]{itemsep=2pt, topsep=4pt}

% --- final-answer box -------------------------------------------------------
% Every problem file ends with its result in one of these. The condensed
% solutions document is assembled by lifting exactly these.
\newtcolorbox{answerbox}[1][]{
  colback=Green!6, colframe=Green!55!black,
  boxrule=0.6pt, arc=2pt, left=6pt, right=6pt, top=4pt, bottom=4pt,
  #1
}

% --- citation to course material --------------------------------------------
% \cite{Module 01, PDF p55}  ->  a small italic parenthetical.
% Solutions cite the SLIDE/PAGE a formula came from so a grader (or a future
% reader) can check the claim against what was actually taught. See the
% directive: every non-obvious step carries one.
\newcommand{\srcref}[1]{{\small\itshape(#1)}}

% --- a step that came from OUTSIDE the course materials ---------------------
% Used only where the course is genuinely silent and the problem asks for
% outside research. Marked visually so it is never mistaken for lecture content.
\newcommand{\extref}[1]{{\small\itshape\color{Sepia}[external: #1]}}

% --- callout for the trap in a problem --------------------------------------
\newtcolorbox{trap}[1][]{
  colback=BurntOrange!7, colframe=BurntOrange!70!black,
  boxrule=0.6pt, arc=2pt, left=6pt, right=6pt, top=4pt, bottom=4pt,
  fonttitle=\bfseries, title=Watch out, #1
}

% --- per-problem header -----------------------------------------------------
% \problemheader{8}{15 pts}{Target clock rate}
% The optional 4-argument form carries a learning outcome where the course
% states one (CESC 410 does; CESC 470's HW1 does not).
\newcommand{\problemheader}[3]{%
  \begin{center}\large\bfseries
    Problem #1 \quad\normalsize\mdseries(#2)\\[2pt]
    \large\bfseries #3
  \end{center}
  \vspace{2pt}\hrule\vspace{8pt}
}
\newcommand{\problemheaderlo}[4]{%
  \begin{center}\large\bfseries
    Problem #1 \quad\normalsize\mdseries(#2, #3)\\[2pt]
    \large\bfseries #4
  \end{center}
  \vspace{2pt}\hrule\vspace{8pt}
}

% Course name in the running head. Defined BEFORE \lhead references it, and
% \renewcommand'd by each course's macros file (inputted after this one).
\providecommand{\coursename}{}

\pagestyle{fancy}
\fancyhf{}
\lhead{\small\coursename}
\rhead{\small\thepage}
\renewcommand{\headrulewidth}{0.3pt}

%   % CESC 470 Computer Architecture — course-specific macros ONLY.
%
% Inputted AFTER docs/latex/coursework_preamble.tex, which carries everything
% shared (answerbox, problemheader, srcref, page style). Put nothing general
% here: if another course would want it, it belongs in the shared preamble.

\renewcommand{\coursename}{CESC 470}

% --- performance-equation notation ------------------------------------------
% The course writes these in words on the slides; these render them compactly
% and, more importantly, IDENTICALLY across every problem file.
\newcommand{\CPUtime}{T_{\mathrm{CPU}}}
\newcommand{\IC}{\mathrm{IC}}                    % instruction count
\newcommand{\CPI}{\mathrm{CPI}}
\newcommand{\cycletime}{t_{\mathrm{cycle}}}
\newcommand{\clockrate}{f_{\mathrm{clk}}}
\newcommand{\cycles}{N_{\mathrm{cycles}}}
\newcommand{\speedup}{S}

% Amdahl's law, in the form the slides state it (Module 01, PDF p86, slide 49).
\newcommand{\amdahl}[3]{#1 + \dfrac{#2}{#3}}

%   \begin{document}
%
% Build-check one file, or a whole assignment:
%   docs/latex/build_tex.sh content/cesc_470/hw/hw01/p01_five_components.tex
%   docs/latex/build_tex.sh content/cesc_470/hw/hw01
%
% Doctrine: docs/directives/coursework-solutions.md
% ─────────────────────────────────────────────────────────────────────────────

\documentclass[11pt]{article}

\usepackage[margin=1in]{geometry}
\usepackage{amsmath, amssymb, mathtools}
\usepackage{graphicx}
\usepackage{float}
\usepackage{booktabs}
\usepackage{longtable}
\usepackage{enumitem}
\usepackage{siunitx}
\usepackage[dvipsnames]{xcolor}
\usepackage{tcolorbox}
\usepackage{fancyhdr}
\usepackage{caption}
\usepackage{tikz}
\usepackage{pgfplots}
\pgfplotsset{compat=1.18}
\usepackage[colorlinks=true, allcolors=NavyBlue]{hyperref}

\captionsetup{font=small, labelfont=bf}
\setlist[enumerate]{itemsep=2pt, topsep=4pt}
\setlist[itemize]{itemsep=2pt, topsep=4pt}

% --- final-answer box -------------------------------------------------------
% Every problem file ends with its result in one of these. The condensed
% solutions document is assembled by lifting exactly these.
\newtcolorbox{answerbox}[1][]{
  colback=Green!6, colframe=Green!55!black,
  boxrule=0.6pt, arc=2pt, left=6pt, right=6pt, top=4pt, bottom=4pt,
  #1
}

% --- citation to course material --------------------------------------------
% \cite{Module 01, PDF p55}  ->  a small italic parenthetical.
% Solutions cite the SLIDE/PAGE a formula came from so a grader (or a future
% reader) can check the claim against what was actually taught. See the
% directive: every non-obvious step carries one.
\newcommand{\srcref}[1]{{\small\itshape(#1)}}

% --- a step that came from OUTSIDE the course materials ---------------------
% Used only where the course is genuinely silent and the problem asks for
% outside research. Marked visually so it is never mistaken for lecture content.
\newcommand{\extref}[1]{{\small\itshape\color{Sepia}[external: #1]}}

% --- callout for the trap in a problem --------------------------------------
\newtcolorbox{trap}[1][]{
  colback=BurntOrange!7, colframe=BurntOrange!70!black,
  boxrule=0.6pt, arc=2pt, left=6pt, right=6pt, top=4pt, bottom=4pt,
  fonttitle=\bfseries, title=Watch out, #1
}

% --- per-problem header -----------------------------------------------------
% \problemheader{8}{15 pts}{Target clock rate}
% The optional 4-argument form carries a learning outcome where the course
% states one (CESC 410 does; CESC 470's HW1 does not).
\newcommand{\problemheader}[3]{%
  \begin{center}\large\bfseries
    Problem #1 \quad\normalsize\mdseries(#2)\\[2pt]
    \large\bfseries #3
  \end{center}
  \vspace{2pt}\hrule\vspace{8pt}
}
\newcommand{\problemheaderlo}[4]{%
  \begin{center}\large\bfseries
    Problem #1 \quad\normalsize\mdseries(#2, #3)\\[2pt]
    \large\bfseries #4
  \end{center}
  \vspace{2pt}\hrule\vspace{8pt}
}

% Course name in the running head. Defined BEFORE \lhead references it, and
% \renewcommand'd by each course's macros file (inputted after this one).
\providecommand{\coursename}{}

\pagestyle{fancy}
\fancyhf{}
\lhead{\small\coursename}
\rhead{\small\thepage}
\renewcommand{\headrulewidth}{0.3pt}

% CESC 470 Computer Architecture — course-specific macros ONLY.
%
% Inputted AFTER docs/latex/coursework_preamble.tex, which carries everything
% shared (answerbox, problemheader, srcref, page style). Put nothing general
% here: if another course would want it, it belongs in the shared preamble.

\renewcommand{\coursename}{CESC 470}

% --- performance-equation notation ------------------------------------------
% The course writes these in words on the slides; these render them compactly
% and, more importantly, IDENTICALLY across every problem file.
\newcommand{\CPUtime}{T_{\mathrm{CPU}}}
\newcommand{\IC}{\mathrm{IC}}                    % instruction count
\newcommand{\CPI}{\mathrm{CPI}}
\newcommand{\cycletime}{t_{\mathrm{cycle}}}
\newcommand{\clockrate}{f_{\mathrm{clk}}}
\newcommand{\cycles}{N_{\mathrm{cycles}}}
\newcommand{\speedup}{S}

% Amdahl's law, in the form the slides state it (Module 01, PDF p86, slide 49).
\newcommand{\amdahl}[3]{#1 + \dfrac{#2}{#3}}


\begin{document}

\problemheader{9}{12 pts --- a: 5, b: 7}{Average CPI and execution time}

\subsection*{Problem statement}

A program has the following instruction mix:

\begin{center}
\begin{tabular}{lll}
  \toprule
  Class & Share of instructions & Cycles per instruction \\
  \midrule
  A & 40\% & 1 \\
  B & 40\% & 2 \\
  C & 20\% & 3 \\
  \bottomrule
\end{tabular}
\end{center}

\begin{enumerate}[label=(\alph*)]
  \item Compute the average CPI. \hfill (5 points)
  \item If the machine's clock cycle time is 1 ns, calculate the total execution
    time for 1 million instructions. \hfill (7 points)
\end{enumerate}

\subsection*{Approach}

This is the course's CPI Example II \srcref{Module 01, PDF p75--76, slides
42--43} in percentage form rather than instruction counts. Average CPI is the
\emph{weighted mean} of the per-class costs:
\begin{equation}
  \CPI_{\text{avg}} = \sum_{i} f_i \times \CPI_i,
  \qquad f_i = \text{fraction of instructions in class } i,
  \label{eq:cpiavg}
\end{equation}
which is the fractional form of the slide's
$\CPI = (\text{total cycles})/(\text{instruction count})$.

Part (b) is then the full performance equation \srcref{PDF p55, slide 32},
in the three-factor form:
\begin{equation}
  \CPUtime = \IC \times \CPI \times \cycletime.
  \label{eq:three}
\end{equation}

Note the fractions are given as percentages, so no instruction count is needed
for (a) --- weights are relative. Part (b) supplies the count.

\subsection*{Work}

\subsubsection*{(a) Average CPI}

Convert percentages to fractions and apply \eqref{eq:cpiavg}:
\begin{align}
  \CPI_{\text{avg}} &= (0.40)(1) + (0.40)(2) + (0.20)(3) \\
                    &= 0.40 + 0.80 + 0.60 \\
                    &= \mathbf{1.8\ \text{cycles/instruction.}}
\end{align}

Sanity on the weights: $0.40 + 0.40 + 0.20 = 1.00$ \checkmark{} --- they must sum
to 1, or the mix is not fully specified.

\paragraph{Equivalent route, per the slide's method.}
The slide works with counts rather than fractions \srcref{PDF p76, slide 43}.
Taking any convenient total, say 10 instructions (4 of A, 4 of B, 2 of C):
\begin{equation}
  \text{cycles} = (4)(1) + (4)(2) + (2)(3) = 4 + 8 + 6 = 18,
  \qquad \CPI = \frac{18}{10} = 1.8 \checkmark
\end{equation}
The answer is independent of the total chosen, which is what makes the
percentage form legitimate.

\subsubsection*{(b) Execution time for $10^{6}$ instructions}

Substitute into \eqref{eq:three} with $\IC = 1{,}000{,}000$,
$\CPI = 1.8$, $\cycletime = 1\ \text{ns} = 1\times10^{-9}\ \text{s}$:
\begin{align}
  \CPUtime &= \IC \times \CPI \times \cycletime \\
           &= (1\times10^{6}) \times 1.8 \times (1\times10^{-9}\ \text{s}) \\
           &= 1.8\times10^{-3}\ \text{s} \\
           &= \mathbf{1.8\ \text{ms.}}
\end{align}

Working it in two visible stages, which is what ``show your work'' wants here:
\begin{align}
  \text{total cycles} &= \IC \times \CPI = 10^{6} \times 1.8
                       = 1.8\times10^{6}\ \text{cycles}, \\
  \CPUtime &= 1.8\times10^{6}\ \text{cycles} \times 1\ \text{ns/cycle}
            = 1.8\times10^{6}\ \text{ns} = 1.8\ \text{ms.}
\end{align}

\begin{trap}
$\cycletime = 1$ ns means the \emph{clock rate is 1 GHz}, not 1 MHz --- the two
are reciprocals \srcref{PDF p54, slide 31}. Inverting this gives an answer off
by $10^{6}$. Also note the answer is 1.8 \emph{milli}seconds: $10^{6}$
instructions at $10^{-9}$ s per cycle is $10^{-3}$ s, not seconds.
\end{trap}

\subsection*{Check}

\paragraph{Independent route for (b), via clock rate.}
$\cycletime = 1$ ns $\Rightarrow \clockrate = 1/(1\times10^{-9}) = 10^{9}$ Hz
$= 1$ GHz. Then, using $\CPUtime = \cycles/\clockrate$:
\begin{equation}
  \CPUtime = \frac{1.8\times10^{6}\ \text{cycles}}{1\times10^{9}\ \text{cycles/s}}
           = 1.8\times10^{-3}\ \text{s} \checkmark
\end{equation}
The same answer without ever multiplying by $\cycletime$ directly.

\paragraph{Bounds check on the CPI.}
The per-class costs run from 1 to 3 cycles, so any weighted average \emph{must}
lie in $[1, 3]$. $1.8$ does. Moreover, since the mix is weighted toward the
cheaper classes (80\% of instructions cost 1 or 2 cycles), the average should sit
below the midpoint of 2 --- and $1.8 < 2$ \checkmark.

\paragraph{Units.}
$$\underbrace{[\text{instructions}]}_{\IC} \times
  \underbrace{\left[\tfrac{\text{cycles}}{\text{instruction}}\right]}_{\CPI} \times
  \underbrace{\left[\tfrac{\text{seconds}}{\text{cycle}}\right]}_{\cycletime}
  = [\text{seconds}] \checkmark$$

\subsection*{Answer}

\begin{answerbox}
\textbf{(a)} $\;\CPI_{\text{avg}} = (0.40)(1) + (0.40)(2) + (0.20)(3)
  = 0.4 + 0.8 + 0.6 = \mathbf{1.8}$ cycles/instruction.

\medskip
\textbf{(b)} $\;\CPUtime = \IC \times \CPI \times \cycletime
  = (10^{6})(1.8)(1\ \text{ns})
  = 1.8\times10^{6}\ \text{ns} = \mathbf{1.8\ \text{ms}}
  \;(= 1.8\times10^{-3}\ \text{s})$.

\medskip
{\small Equivalently: total cycles $= 10^{6} \times 1.8 = 1.8\times10^{6}$; at
1 ns per cycle (a 1 GHz clock) that is 1.8 ms.}
\end{answerbox}

\end{document}
